\section{Perturbation and a generating limit}
\label{sec:generation}

We combine the flow estimate with the small-cocycle construction to
approximate the extra generator by a word in a slightly perturbed
$n$-tuple.  We then arrange these approximations so that the limiting
$n$-tuple generates the entire von Neumann algebra.
We call $C$ a \emph{completion} of $A=(A_1,\ldots,A_n)$ when
$(A,C)$ is a freely generating Haar tuple in $M$.

\subsection{One-step perturbation}

\begin{proposition}[Approximating the extra generator in one step]
\label{prop:perturbation}
Let $n\geq 3$, and let $(A_1,\ldots,A_n,C)$ be a freely generating Haar
tuple in a finite von Neumann algebra $(M,\tau)$.  For every
$\varepsilon>0$ there are a trace-preserving normal automorphism
$\alpha$ of $M$ and a word $w\in\F_n$ such that, with
$A_j'=\alpha(A_j)$,
\begin{equation}
\label{eq:perturbation}
 \max_{1\leq j\leq n}\norm{A_j'-A_j}<\varepsilon,
 \qquad
 \norm{A_w'-C}<\varepsilon.
\end{equation}
The tuple $(A_1',\ldots,A_n',C')$, where $C'=\alpha(C)$, is again freely
generating Haar.  Here the target in the second estimate is the
original $C$; the new completion is
\[
 C'=C A_w^*,
\]
where $A_w$ is evaluated on the original $n$-tuple.
\end{proposition}

\begin{proof}
Choose $0<\eta<\varepsilon/K$, with $K=3\pi$ as in
\eqref{eq:free-sum}.  Lemma~\ref{lem:small-cocycle} gives a finitely
supported real coefficient vector $h_1$, with $h_j=0$ for $j\geq2$,
and a word $w\in\F_n$ for which
\[
 \|h_1\|_{\ell^2}<\eta,
 \qquad
 \|D_w-\e\|_{\ell^2}<\eta.
\]
Apply Proposition~\ref{prop:flow} to these coefficients and write
$\beta=\beta_1$.  Its displacement estimates
\eqref{eq:generator-displacement} and \eqref{eq:word-displacement}
give
\[
 \max_j\norm{\beta(A_j)-A_j}<\varepsilon,
 \qquad
 \norm{\beta(A_w)-C A_w}<\varepsilon.
\]
Moreover, $\beta$ fixes $S$ and therefore fixes $C=\exp(\ii S)$.

The free-group basis transformation that fixes the first $n$ symbols
and sends the last symbol $c$ to $c w$ is invertible: its inverse
fixes the first $n$ symbols and sends $c$ to $c w^{-1}$.  By
Lemma~\ref{lem:haar-identification}, it induces a trace-preserving
normal automorphism $\gamma$ of $M$ satisfying
\[
 \gamma(A_j)=A_j \quad(1\leq j\leq n),
 \qquad
 \gamma(C)=C A_w.
\]
Set $\alpha=\gamma^{-1}\circ\beta$.  Since $\gamma$ fixes every
$A_j$ and hence $A_w$, we have the exact identities
\begin{align*}
 \alpha(A_j)-A_j
   &=\gamma^{-1}\bigl(\beta(A_j)-A_j\bigr),\\
 \alpha(A_w)-C
   &=\gamma^{-1}\bigl(\beta(A_w)-C A_w\bigr).
\end{align*}
An automorphism is isometric, so these identities prove
\eqref{eq:perturbation}, with the displayed order of the factors.
The image under $\alpha$ of the original freely generating Haar
tuple is again such a tuple.  Finally,
\[
 C'=\alpha(C)=\gamma^{-1}(C)=C A_w^*,
\]
because $\beta(C)=C$ and
$\gamma(C A_w^*)=C A_w A_w^*=C$.
\end{proof}

\subsection{A generating limit}

We record the elementary continuity estimate that will control the
iteration.  For two $d$-tuples of unitaries $U=(U_1,\ldots,U_d)$ and
$V=(V_1,\ldots,V_d)$, put
\[
 d_2(U,V)=\max_{1\leq j\leq d}\normtwo{U_j-V_j}.
\]
If $v$ is a group word of length $|v|$, telescoping its product gives
\[
 \normtwo{U_v-V_v}\leq |v|\,d_2(U,V).
\]
Indeed, multiplication on either side by a unitary preserves the
tracial $L^2$ norm, and
$\normtwo{U_j^*-V_j^*}=\normtwo{U_j-V_j}$.
Consequently, a finite word polynomial
$F=\sum_{\nu=1}^q c_\nu v_\nu$ satisfies
\begin{equation}
\label{eq:word-polynomial-lipschitz}
 \normtwo{F(U)-F(V)}
 \leq L(F)d_2(U,V),
 \qquad
 L(F)=\sum_{\nu=1}^q |c_\nu|\,|v_\nu|.
\end{equation}
The same telescoping argument also gives continuity in operator norm.

\begin{proof}[Proof of Theorem~\ref{thm:rank-step}]
Fix $n\geq3$ and let $M=L(\F_{n+1})$, with its canonical normalized
trace and canonical freely generating Haar tuple
\[
 (A_1^{(0)},\ldots,A_n^{(0)},C^{(0)}).
\]
Write $A^{(k)}=(A_1^{(k)},\ldots,A_n^{(k)})$.  Choose a sequence
$(y_j)_{j\geq1}$ that is dense in $L^2(M,\tau)$ and consists of
elements of the rational complex group word algebra of the initial
tuple.  Such a sequence exists by
Lemma~\ref{lem:haar-identification}; coefficients in
$\mathbb{Q}+\ii\mathbb{Q}$ make this algebra countable.  Put $r_0=1$.

We construct freely generating Haar tuples $(A^{(k)},C^{(k)})$ and
positive numbers $r_k$ by induction.  Suppose the construction has
reached stage $k-1$, where $k\geq1$.  The choices at stage $k$ are
made in the following order.

First, for each $1\leq j\leq k$, choose a finite linear combination
$F_{j,k}$ of group words in $n+1$ variables such that
\begin{equation}
\label{eq:stage-old-approximation}
 \normtwo{
 F_{j,k}(A^{(k-1)},C^{(k-1)})-y_j
 }<2^{-k}.
\end{equation}
This is possible because the current completed tuple generates $M$
and its group word algebra is dense in $L^2(M,\tau)$.
Let $H_k\geq0$ be a common bound for the finitely many constants
$L(F_{j,k})$ in \eqref{eq:word-polynomial-lipschitz}.  Choose
\begin{equation}
\label{eq:stage-tolerance}
 0<\varepsilon_k<
 \min\left\{\frac{r_{k-1}}2,\,
                  \frac{2^{-k}}{1+H_k}\right\}.
\end{equation}

Next apply Proposition~\ref{prop:perturbation} to
$(A^{(k-1)},C^{(k-1)})$ with tolerance $\varepsilon_k$.
Let $\alpha_k$ and $w_k\in\F_n$ be its automorphism and word, and set
\[
 A_j^{(k)}=\alpha_k(A_j^{(k-1)}),
 \qquad
 C^{(k)}=\alpha_k(C^{(k-1)}).
\]
The new completed tuple is freely generating Haar, and
\begin{equation}
\label{eq:stage-perturbation}
 \max_j\norm{A_j^{(k)}-A_j^{(k-1)}}<\varepsilon_k,
 \qquad
 \norm{A_{w_k}^{(k)}-C^{(k-1)}}<\varepsilon_k.
\end{equation}
In particular, the last expression is close to the old completion,
as needed in \eqref{eq:stage-old-approximation}.

Now define word polynomials in $n$ variables by
\[
 G_{j,k}(X)=F_{j,k}(X,X_{w_k})
 \qquad(1\leq j\leq k).
\]
Since $\normtwo{x}\leq\norm{x}$, equations
\eqref{eq:word-polynomial-lipschitz},
\eqref{eq:stage-old-approximation}, and
\eqref{eq:stage-perturbation} imply
\begin{align}
\label{eq:stage-new-approximation}
 \normtwo{G_{j,k}(A^{(k)})-y_j}
 &\leq
 \normtwo{F_{j,k}(A^{(k)},A_{w_k}^{(k)})
       -F_{j,k}(A^{(k-1)},C^{(k-1)})}
 \notag\\
 &\qquad
 +\normtwo{F_{j,k}(A^{(k-1)},C^{(k-1)})-y_j}
 \notag\\
 &<H_k\varepsilon_k+2^{-k}<2\cdot2^{-k}.
\end{align}
This estimate uses $A_{w_k}^{(k)}$ as a single unitary slot of
$F_{j,k}$.  Thus the tolerance chosen in
\eqref{eq:stage-tolerance} does not depend on the length of the word
that the perturbation subsequently produces.

Finally, let $L_k\geq0$ be a common $L^2$-Lipschitz bound for the
finitely many polynomials $G_{j,k}$, and choose
\begin{equation}
\label{eq:future-budget}
 0<r_k\leq
 \min\left\{\frac{r_{k-1}}2,\,
                  \frac{2^{-k}}{1+L_k}\right\}.
\end{equation}
These constants are finite by
\eqref{eq:word-polynomial-lipschitz}.  They are chosen after $w_k$
is known, so any increase in word length is accommodated in the
budget for future steps.  Both quantities in the minimum are
positive, and the induction can continue.

Figure~\ref{fig:iteration} records this order of choices and the
role of the budget in controlling all later perturbations.
\begin{figure}[htbp]
\centering
\begin{tikzpicture}[
  stage/.style={draw=blue!45!black, rounded corners=2pt,
    fill=blue!3, text width=3.65cm, minimum height=2.1cm,
    align=center, inner sep=7pt, font=\small},
  choice/.style={-{Stealth[length=2.3mm]}, thick, draw=blue!45!black},
  node distance=9mm and 9mm
]
\node[stage] (old) {$F_{j,k}$\\[3pt]
  Approximate $y_1,\ldots,y_k$\\in the current tuple};
\node[stage, right=of old] (epsilon) {$\varepsilon_k$\\[3pt]
  Choose the tolerance\\from $H_k$ and $r_{k-1}$};
\node[stage, right=of epsilon] (word) {$\alpha_k,\ w_k$\\[3pt]
  Perturb the tuple;\\approximate the old $C^{(k-1)}$};
\node[stage, below=of word] (new) {$G_{j,k}$\\[3pt]
  Substitute the witness word:\\
  $F_{j,k}(X,X_{w_k})$};
\node[stage, below=of epsilon] (budget) {$L_k,\ r_k$\\[3pt]
  Choose a\\Lipschitz bound;\\fix the future budget};
\node[stage, below=of old] (tail) {Tail control\\[3pt]
  $\sum_{\ell>k}\varepsilon_\ell\leq r_k$\\
  Preserve earlier approximations};
\draw[choice] (old) -- (epsilon);
\draw[choice] (epsilon) -- (word);
\draw[choice] (word) -- (new);
\draw[choice] (new) -- (budget);
\draw[choice] (budget) -- (tail);
\end{tikzpicture}

\caption{The choices at stage $k$ follow the arrows.  The witness
word $w_k$ is found only after $\varepsilon_k$ is fixed.  A longer
witness word may increase the Lipschitz bound $L_k$ of the substituted
polynomials $G_{j,k}$, so its length affects the future budget $r_k$
only.  This budget bounds the total change after stage $k$ and
preserves the target approximations in the limit.}
\label{fig:iteration}
\end{figure}

For $\ell>k$, the geometric decrease of the budgets gives
\[
 \varepsilon_\ell<\frac{r_{\ell-1}}2
       \leq \frac{r_k}{2^{\ell-k}}.
\]
Using \eqref{eq:stage-perturbation}, we therefore obtain, for $p>k$,
\[
 \max_j\norm{A_j^{(p)}-A_j^{(k)}}
 \leq\sum_{\ell=k+1}^{p}\varepsilon_\ell
 \leq r_k.
\]
Since $r_k\leq2^{-k}r_0$, each sequence $(A_j^{(k)})_k$ converges in
operator norm to a unitary $A_j^{(\infty)}$, and
\begin{equation}
\label{eq:generator-tail}
 \max_j\norm{A_j^{(\infty)}-A_j^{(k)}}\leq r_k.
\end{equation}
The limits are unitaries because the unitary group is norm closed.
For every nonidentity reduced word $v\in\F_n$,
\[
 \tau(A_v^{(\infty)})
   =\lim_{k\to\infty}\tau(A_v^{(k)})=0,
\]
by operator-norm continuity of each fixed word evaluation and of
the trace.  Thus the limiting tuple has the free Haar word
distribution.

It remains to prove that this tuple generates $M$.  For every
$j\leq k$, equations \eqref{eq:stage-new-approximation},
\eqref{eq:future-budget}, and \eqref{eq:generator-tail} yield
\begin{align}
\label{eq:limit-target-approximation}
 \normtwo{G_{j,k}(A^{(\infty)})-y_j}
 &\leq
 \normtwo{G_{j,k}(A^{(\infty)})-G_{j,k}(A^{(k)})}
  +\normtwo{G_{j,k}(A^{(k)})-y_j}
 \notag\\
 &\leq L_k r_k+2\cdot2^{-k}
 \leq3\cdot2^{-k}.
\end{align}
Fixing $j$ and letting $k\to\infty$ through $k\geq j$ shows that
$y_j$ belongs to the $L^2$ closure of the word algebra of
$A^{(\infty)}$.  Since the $y_j$ are dense, this word algebra is
dense in all of $L^2(M,\tau)$.

To pass from this density to von Neumann generation, set
$N=W^*(A_1^{(\infty)},\ldots,A_n^{(\infty)})$.  The standard
finite-trace conditional expectation $E_N:M\to N$ is
trace-preserving and extends to the orthogonal projection of
$L^2(M,\tau)$ onto $L^2(N,\tau)$.  In particular it is
$L^2$-contractive and fixes every element of $N$.
Each $G_{j,k}(A^{(\infty)})$ belongs to $N$, so
\[
 \normtwo{E_N(y_j)-y_j}
 \leq 2\normtwo{G_{j,k}(A^{(\infty)})-y_j}
 \longrightarrow0.
\]
Hence the projection fixes the dense family $(y_j)$ and is the
identity on $L^2(M,\tau)$.  For $x\in M$, this implies
$E_N(x)=x$ by faithfulness of $\tau$, and therefore $N=M$.

The limiting tuple is now a freely generating Haar $n$-tuple in
$M$.  Lemma~\ref{lem:haar-identification} supplies a unital
trace-preserving normal $*$-isomorphism from $L(\F_n)$ onto
$M=L(\F_{n+1})$, proving the theorem.
\end{proof}

\begin{remark}
Only the first $n$ generators are required to converge in operator
norm.  The completions $C^{(k)}$ may vary without a limit.
The target approximations are in $L^2$, and the conditional
expectation argument permits the operator norms of the
approximating polynomials to grow with $k$.
\end{remark}
